\subsection{卷积与基本解}\label{sec:convolutions-fundamental-solutions}

定义在 $\mathbb{R}^m$ 上的两个函数的卷积的经典定义为
\begin{equation}\label{eq:classical-convolution-definition}
f*g(\bm{x})=\int_{\mathbb{R}^m}f(\bm{x}-\bm{y})g(\bm{y})\,\mathrm{d}\bm{y}.
\end{equation}
本节将考虑把这一定义推广到广义函数, 并给出其在求解常系数偏微分方程中的应用.

\subsubsection{分布的直积}\label{subsec:distribution-direct-product}

一般不能定义两个广义函数 $f(\bm{x})$ 和 $g(\bm{x})$ 的乘积. 然而, 总可以将依赖于不同变量的两个广义函数相乘. 即若 $f\in\mathcal{D}^{\prime}(\mathbb{R}^p)$ 且 $g\in\mathcal{D}^{\prime}(\mathbb{R}^q)$, 则可以将 $f(\bm{x})g(\bm{y})$ 定义为 $\mathbb{R}^{p+q}$ 上的分布.

\begin{Definition}
\label{def:distribution-direct-product}
设 $f\in\mathcal{D}^{\prime}(\mathbb{R}^p)$, $g\in\mathcal{D}^{\prime}(\mathbb{R}^q)$. 那么\textbf{直积} $f(\bm{x})g(\bm{y})$ 是由下式给出的 $\mathbb{R}^{p+q}$ 上的分布
\begin{equation}\label{eq:distribution-direct-product-definition}
(f(\bm{x})g(\bm{y}),\phi(\bm{x},\bm{y}))=(f(\bm{x}),(g(\bm{y}),\phi(\bm{x},\bm{y}))).
\end{equation}

\end{Definition}

即首先将 $\phi(\bm{x},\bm{y})$ 视为仅关于 $\bm{y}$ 的函数, 它以 $\bm{x}$ 为参数. 对这一函数应用泛函 $g$. 所得结果为实值函数 $\psi(\bm{x})$, 它显然具有紧支集. 容易证明 $\psi$ 属于 $C^\infty$ 类 (Problem~\ref{ex:distribution-smooth-test-parameter-map}). 因而 $\psi$ 为试验函数, $(f,\psi)$ 有良好定义. 为说明这一定义合理, 还需证明: 若 $\phi_n$ 在 $\mathcal{D}(\mathbb{R}^{p+q})$ 中收敛到零, 则 $(f(\bm{x}),(g(\bm{y}),\phi_n(\bm{x},\bm{y})))$ 收敛到零. 由于 $f$ 为分布, 只需证明 $\psi_n:=(g(\bm{y}),\phi_n(\bm{x},\bm{y}))$ 在 $\mathcal{D}(\mathbb{R}^p)$ 中收敛到零. 若 $S_p\times S_q$ 为包含所有 $\phi_n$ 的支集的紧集, 则 $S_p$ 包含所有 $\psi_n$ 的支集. 剩下需要证明 $\psi_n$ 及其所有导数一致收敛到零. 设 $\alpha$ 为 $p$ 维多重指标, 令 $\beta=(\alpha,0,\ldots,0)$. 假设 $D^\alpha\psi_n$ 不一致收敛到零. 选取子列后, 可假设存在点列 $\bm{x}_n$, 使
\begin{equation}\label{eq:direct-product-continuity-contradiction}
|D^\alpha\psi_n(\bm{x}_n)|=|(g(\bm{y}),D^\beta\phi_n(\bm{x}_n,\bm{y}))|\geq\epsilon.
\end{equation}
但由于 $\phi_n$ 连同所有导数均收敛到零, 序列 $\chi_n(\bm{y}):=D^\beta\phi_n(\bm{x}_n,\bm{y})$ 也如此. 因而 $\chi_n$ 在 $\mathcal{D}(\mathbb{R}^q)$ 中收敛到零, 所以 $(g,\chi_n)$ 收敛到零, 矛盾.

\begin{Example}
\label{ex:direct-product-delta}
作为直积的一个简单例子, 注意到 \[\delta(\bm{x})\delta(\bm{y})=\delta(\bm{x},\bm{y}).\]
\end{Example}

若 $\phi(\bm{x},\bm{y})$ 具有特殊形式 $\phi_1(\bm{x})\phi_2(\bm{y})$, 则得到
\begin{equation}\label{eq:direct-product-separable-test-pairing}
(f(\bm{x})g(\bm{y}),\phi(\bm{x},\bm{y}))=(f,\phi_1)(g,\phi_2).
\end{equation}


形如 $\phi_1(\bm{x})\phi_2(\bm{y})$ 的函数的线性组合实际上在 $\mathcal{D}(\mathbb{R}^{p+q})$ 中稠密. 为说明这一点, 设 $\phi$ 的支集位于集合 $Q:=\{|\bm{x}|\leq a,\ |\bm{y}|\leq a\}$ 内. 由 Weierstrass 逼近定理 (见 Section~\ref{subsec:weierstrass-approximation}), 存在多项式序列, 在集合 $Q^{\prime}:=\{|\bm{x}|\leq2a,\ |\bm{y}|\leq2a\}$ 上一致收敛到 $\phi$. 此外, 定理证明中的论证还表明, 多项式的导数在 $Q^{\prime}$ 上一致收敛到 $\phi$ 的相应导数. 因此可以选取多项式 $p_n$, 使在 $Q^{\prime}$ 上有
\begin{equation}\label{eq:direct-product-polynomial-derivative-approximation}
|D^\alpha p_n-D^\alpha\phi|\leq\dfrac1n,\quad\forall\,|\alpha|\leq n.
\end{equation}
现在设 $b_1(\bm{x})$, $b_2(\bm{y})$ 为固定的试验函数, 当 $|\bm{x}|\leq a$ ($|\bm{y}|\leq a$) 时等于 1, 当 $|\bm{x}|\geq2a$ ($|\bm{y}|\geq2a$) 时等于 0. 那么序列
\begin{equation}\label{eq:direct-product-cutoff-polynomial-approximation}
\phi_n(\bm{x},\bm{y}):=b_1(\bm{x})b_2(\bm{x})p_n(\bm{x},\bm{y})
\end{equation}
在 $\mathcal{D}(\mathbb{R}^{p+q})$ 中收敛到 $\phi$.

利用这一事实与连续性, 只需在形如 $\phi_1(\bm{x})\phi_2(\bm{y})$ 的受限试验函数集合上验证, 就能证明直积的性质. 一个直接结论是, 在定义中可以按相反次序计算 $f$ 与 $g$, 即还有
\begin{equation}\label{eq:direct-product-reversed-pairing}
(f(\bm{x})g(\bm{y}),\phi(\bm{x},\bm{y}))=(g(\bm{y}),(f(\bm{x}),\phi(\bm{x},\bm{y})));
\end{equation}
用如下具有启示性的记号表示这一事实
\begin{equation}\label{eq:distribution-direct-product-commutativity}
f(\bm{x})g(\bm{y})=g(\bm{y})f(\bm{x}).
\end{equation}


另一个明显的性质是结合律
\begin{equation}\label{eq:distribution-direct-product-associativity}
f(\bm{x})(g(\bm{y})h(\bm{z}))=(f(\bm{x})g(\bm{y}))h(\bm{z}).
\end{equation}

